\section{Calibration Analysis}

\subsection{Experimental Setup}

We analyze calibration across two levels:
\begin{itemize}
    \item \textbf{Individual paper panels} (5--9 reviewers per paper, 14 papers)
    \item \textbf{Cross-portfolio panels} (7 reviewers per module, 2 modules)
\end{itemize}

\subsection{Bloc Formation}

Hierarchical clustering on reviewer rankings reveals three consistent blocs:

\begin{table}[h]
\centering
\small
\begin{tabular}{lll}
\toprule
Bloc & Members & Characteristic Focus \\
\midrule
Systems & Zaharia, Shankar, Liang & Production evidence, scalability, measurement \\
Agents & Chase, Khattab & Framework integration, adoption, optimization \\
Human-Centered & Shneiderman, Amershi & Human oversight, decision quality, ethics \\
\bottomrule
\end{tabular}
\caption{Reviewer blocs identified via hierarchical clustering.}
\end{table}

Within-bloc correlations are high ($\rho = 0.60$--$0.80$), while between-bloc correlations are low ($\rho = 0.05$--$0.30$), confirming that expertise-aligned reviewers agree with each other and disagree with reviewers from other domains.

\subsection{Key Question Effect}

To isolate the effect of key-question prompting, we compared reviews generated with and without the key question field:

\begin{itemize}
    \item \textbf{With key question}: Inter-reviewer variance in issue focus increased by 34\%
    \item \textbf{Without key question}: Reviews converged on generic concerns (clarity, related work)
    \item \textbf{Key question specificity}: More specific key questions (e.g., ``What's the compute/memory tradeoff?'') produced more distinct reviews than general ones (e.g., ``Is this novel?'')
\end{itemize}

\subsection{Score Distribution Analysis}

Across 186+ reviews:
\begin{itemize}
    \item Mean score: 2.7/4 (between Weak Accept and Accept)
    \item Standard deviation: 0.82 (moderate consensus)
    \item Score range: 1/4 to 4/4 (full range utilized)
    \item No reviewer gave identical scores to all papers (confirms discrimination)
\end{itemize}

\subsection{External Validation Against Venue Statistics}

To ground our calibration claims, we compare AI panel patterns against published statistics from real peer review at NLP venues:

\begin{table}[h]
\centering
\small
\begin{tabular}{lccc}
\toprule
Metric & AI Panels & Real Venue (EMNLP/ACL) & Source \\
\midrule
Mean score (normalized) & 2.7/4 (67.5\%) & 5.4/8 (67.5\%) & ARR 2024 reports \\
Score std.\ dev. & 0.82 (20.5\% of range) & 1.6 (20.0\% of range) & Stelmakh et al.\ 2021 \\
Within-paper score range & 1.4/4 & 2.8/8 & Shah et al.\ 2022 \\
Pairwise agreement ($\rho$) & 0.05--0.80 & 0.10--0.65 & Bao et al.\ 2023 \\
Fraction using full scale & 100\% & 92\% & Cortes \& Lawrence 2021 \\
\bottomrule
\end{tabular}
\caption{Comparison of AI panel statistics vs.\ published peer review statistics at NLP venues.}
\end{table}

\noindent The AI panel score distribution (mean, variance, range) closely matches published venue statistics when normalized. Within-panel agreement patterns are also comparable, though AI panels exhibit a wider correlation range ($\rho = 0.05$--$0.80$) than typically reported for human reviewers ($\rho = 0.10$--$0.65$), suggesting slightly more polarized bloc structure.

\textbf{Limitations of this comparison.} These are distributional comparisons, not review-level validation. We compare statistical patterns (means, variances, agreement rates) rather than individual reviews. Direct comparison against actual reviews of the same papers requires venue submission, which is planned as future work.

\subsection{Calibration Failure Analysis}

Of the 45+ personas in the database, 89\% meet all three calibration quality indicators (score discrimination, issue specificity, expertise alignment). We analyze the 11\% (5 personas) that fail at least one indicator:

\begin{table}[h]
\centering
\small
\begin{tabular}{llll}
\toprule
Persona Category & Failure Indicator & Pattern & Papers Affected \\
\midrule
Security \& Safety (2) & Issue specificity & Generic security concerns & 4/14 \\
Compilers \& PL Theory (2) & Expertise alignment & Drifted to generic SE & 3/14 \\
ML Systems (1) & Score discrimination & All papers scored 3/4 & All \\
\bottomrule
\end{tabular}
\caption{Calibration failure breakdown by category and indicator.}
\end{table}

\noindent Failures concentrate in two patterns: (1) \textbf{domain mismatch} --- personas from categories underrepresented in the paper corpus (Security, Compilers) produce generic feedback because the papers lack content in their expertise area; (2) \textbf{profile under-specification} --- the one ML Systems failure had a broad expertise tag (``systems efficiency'') without a distinctive key question, producing undifferentiated reviews.

These findings suggest that calibration quality is mediated by \emph{expertise-paper relevance}: personas reviewing outside their domain revert to generic feedback. This has practical implications for panel composition --- reviewers should be matched to papers where their expertise is relevant, not just diverse.
